% Chapter 8 -- Claim: Value Beats Growth
% Requires: amsmath, amssymb, booktabs, array, tabularx, graphicx, tcolorbox (breakable, skins), hyperref or url.
% Defines the callout box only if the preamble has not.
\providecolor{fenavy}{HTML}{1B2A6B}
\providecolor{feshade}{HTML}{F3F5FA}
\makeatletter
\@ifundefined{callout}{%
  \newtcolorbox{callout}{breakable, enhanced,
    colback=white, colframe=gray!70, boxrule=0.5pt, arc=1mm,
    borderline west={2.5pt}{0pt}{gray!70},
    left=8pt, right=8pt, top=4pt, bottom=4pt}%
}{}
\makeatother

\chapter{Claim: Value Beats Growth}
\label{ch:08}

Buying what is cheap is the oldest idea in investing. Benjamin Graham taught it in the 1930s (Graham and Dodd 1934), Warren Buffett built a fortune on a refined version of it, and in the early 1990s academic finance gave it a name, a number and a factor: the value premium\index{value premium}, measured by HML\index{factor!HML (value)}. For three decades afterward, a tilt toward value stocks\index{value versus growth} was one of the most widely recommended departures from the market portfolio, built into index funds, smart-beta\index{beta} products and the asset allocations of large institutions.

Then, from about 2007, value stopped working in the United States. For more than a decade expensive growth stocks beat cheap value stocks by a wide margin, and the debate turned from why value works to whether it ever did, whether it had been arbitraged away, or whether accounting had changed what ``cheap'' means.

This chapter puts the claim on trial\index{claims-on-trial method} with a century of U.S. data, other measures of cheapness, other countries and the investor's own question: even if the premium is real, could you have held on long enough to collect it?

By the end of the chapter you will be able to:

\begin{itemize}
  \item explain the motivation behind value investing research and its implications for factor investing, risk theory, and capital allocation;
  \item distinguish Chapter~\ref{ch:08}'s empirical factor trial from Chapter~\ref{ch:03}'s econometric factor modeling framework;
  \item identify and apply the evaluation tools drawn from Chapters~\ref{ch:01} to~\ref{ch:06} specifically for testing cross-sectional equity anomalies;
  \item compute book-to-market\index{book-to-market ratio} and other valuation ratios, and explain why intangible\index{intangible capital} investment distorts book value;
  \item explain what an anomaly\index{anomaly} is, and how portfolio sorts\index{portfolio sort} and long--short factors are built;
  \item construct HML from the six size and book-to-market portfolios;
  \item compare the risk and mispricing\index{mispricing} explanations of the value premium;
  \item state the claim that value beats growth in a form that can fail, and test it over time, across measures, among large stocks, abroad and after costs;
  \item size a value tilt you could hold through a decade of underperformance.
\end{itemize}

\section{Motivation: why this claim matters and its implications}
\label{sec:08-motivation}

The value premium is the cornerstone of empirical factor pricing and quantitative equity management. Whether cheap stocks systematically outperform expensive ones goes beyond selecting individual stocks; it tests how modern financial markets price business assets and allocate economic capital.

\subsection{Why researchers and practitioners care}

\begin{itemize}
  \item \textbf{Institutional factor allocation:} Value is the foundational building block for smart-beta products, quantitative equity funds, and institutional factor tilts holding trillions of dollars globally.
  \item \textbf{The risk versus mispricing debate:} For academic finance, the value premium is the central arena for testing market efficiency. If value earns a premium because it is riskier in economic downturns, market efficiency holds. If it earns a premium because investors systematically over-extrapolate growth trends, markets are behaviorally inefficient.
  \item \textbf{Career and business survival:} The severe 2007--2020 value drawdown created acute risk for institutional asset managers, forcing a fundamental question: has the value strategy been arbitraged away, is raw accounting data failing in a digital economy, or is the premium merely experiencing a temporary statistical drought?
\end{itemize}

\subsection{Key implications}

The empirical status of the value premium drives decisions across multiple dimensions of finance:

\begin{itemize}
  \item \textbf{For investors and portfolio designers:} If the value premium is real but subject to multi-year droughts, sizing a value tilt requires strict risk budgeting to avoid abandoning the strategy at market troughs.
  \item \textbf{For corporate valuation and cost of capital:} If expensive growth companies systematically earn lower expected returns, their cost of equity capital is lower. Firms with high market-to-book ratios can fund long-duration projects at lower hurdle rates.
  \item \textbf{For quantitative model design:} As corporate investment shifts from physical capital (factories and machinery) to intangible assets (software, R\&D, and brand equity), traditional valuation ratios require structural adjustments to avoid misclassifying innovative firms.
\end{itemize}

\section{Toolkit needed from earlier chapters}
\label{sec:08-connection}

Testing cross-sectional factor claims requires rigorous empirical hygiene to prevent look-ahead bias, data mining, and unrealistic execution assumptions. Understanding how Chapter~\ref{ch:08} connects to earlier methodological chapters clarifies its position in the book.

\subsection{Relationship to Multi-Factor Models (Chapter~\ref{ch:03})}

Because both chapters examine multi-factor models and use Fama--French data, it is essential to distinguish their econometric roles and empirical scopes:

\begin{itemize}
  \item \textbf{Methodology vs.\ empirical trial:} Chapter~\ref{ch:03} is an econometric tool chapter. It establishes the mathematical framework of multi-factor regressions ($R_{i,t} - R_{f,t} = \alpha_i + \beta_i MKT_t + s_i SMB_t + h_i HML_t + \varepsilon_{i,t}$), linear variance partitioning, and robust standard error selection (Newey--West, GMM, bootstrap). Chapter~\ref{ch:08} is a claim chapter. It takes the HML factor as a specific financial claim and evaluates whether cheap stocks beat expensive stocks out of sample and after transaction costs.
  \item \textbf{Lab 3 vs.\ Lab 8 distinction:} Lab 3 teaches econometric diagnostics—how to estimate alphas, partition variance, and evaluate why naive OLS $t$-statistics fail on heavy-tailed financial data. Lab~\ref{ch:08} replicates Fama and French (2021) to put the HML factor itself on trial—testing its performance across four subperiods (1926--1963, 1963--1990, 1991--2006, 2007--present), evaluating alternate cheapness ratios ($E/P$, $C/P$, $D/P$), examining small vs.\ large-cap splits, testing international markets, and sizing a value tilt under drawdown constraints.
\end{itemize}

\subsection{Summary of tools drawn from Chapters 1--6}

Table~\ref{tab:08-tools} explicitly identifies the specific analytical tools drawn from Chapters~\ref{ch:01} through~\ref{ch:06} that are deployed throughout this chapter to evaluate the value premium.

\begin{table}[htbp]
\centering
\small
\caption{Tools drawn from Chapters 1--6 used in Chapter 8.}
\label{tab:08-tools}
\begin{tabularx}{0.95\linewidth}{@{}>{\raggedright\arraybackslash}p{0.18\linewidth}>{\raggedright\arraybackslash}p{0.28\linewidth}>{\raggedright\arraybackslash}X@{}}
\toprule
Origin Chapter & Tool / Concept & Role in Chapter 8 \\
\midrule
Chapter~\ref{ch:01} & Expected vs.\ Realized Returns & Evaluating whether value underperformance reflects lower expected returns or unforecastable shocks. \\
Chapter~\ref{ch:02} & Point-in-Time Data \& Publication Lags & Aligning accounting book equity with market value using 6-month lags (June rebalancing) to prevent look-ahead leakage. \\
Chapter~\ref{ch:03} & CAPM Alpha, Betas \& Robust Standard Errors & Estimating HML's excess return above market risk and using Newey--West/bootstrap corrections for factor significance. \\
Chapter~\ref{ch:04} & Search Hazards \& Multiple Testing & Testing multiple valuation metrics (B/M, E/P, C/P, D/P) across global regions without false discovery. \\
Chapter~\ref{ch:05} & Subperiod Stability \& Out-of-Sample Tests & Comparing the original 1963--1990 sample against the post-2007 drought sample to test factor decay. \\
Chapter~\ref{ch:06} & Transaction Costs \& Turnover Drag & Accounting for HML's 3\% monthly turnover and calculating net performance after execution costs. \\
\bottomrule
\end{tabularx}
\end{table}

\section{Foundations: accounting value, anomalies and factor portfolios}
\label{sec:08-foundations}

\subsection{Book value and market value}

A firm's \textbf{book value of equity}\index{book-to-market ratio|textbf} is an accounting figure from its \textbf{balance sheet}, the statement of what a firm owns and owes: total assets minus total liabilities, built mainly from the historical cost of what the firm has invested, less \textbf{depreciation}, the gradual write-down of assets as they wear out. Its \textbf{market value of equity} is the share price times shares outstanding: what investors think the firm's future cash flows are worth today. The ratio of the two, \textbf{book-to-market} (or its inverse, price-to-book), compares what the firm has put in with what the market expects it to earn.

Accounting rules record spending on factories and equipment as assets, but treat most spending on research, software, brands and training as expenses that reduce earnings and never appear on the balance sheet. Firms that invest mainly in such \textbf{intangible}\index{intangible capital|textbf} assets therefore tend to have low book values relative to their market values (Figure~\ref{fig:08-book}).

\subsection{Other valuation ratios}

Each of these measures how much the investor gets per dollar paid:

\begin{itemize}
  \item \textbf{earnings-to-price}\index{valuation ratios!earnings yield|textbf} (the inverse of the P/E ratio) compares annual profits with the share price;
  \item \textbf{cash flow-to-price}\index{valuation ratios!cash flow-to-price|textbf} uses earnings plus depreciation and other non-cash charges;
  \item \textbf{dividend yield}\index{valuation ratios!dividend yield|textbf} uses dividends paid.
\end{itemize}

Stocks high on these ratios are called \textbf{value} stocks, those low on them \textbf{growth} stocks.

\subsection{The CAPM and its failures}

Chapter~\ref{ch:03} introduced the capital asset pricing model\index{capital asset pricing model (CAPM)} and alpha\index{alpha}, the part of a stock's average return that its market beta does not explain (Section~\ref{sec:03-capm}). From the late 1970s, researchers found that stocks with high earnings-to-price or book-to-market ratios had positive alphas: they earned more than their betas justified. These became known as \textbf{anomalies}\index{anomaly|textbf}, patterns the CAPM could not explain.

\subsection{Portfolio sorts}

The standard empirical method for studying a characteristic like book-to-market is a \textbf{portfolio sort}\index{portfolio sort|textbf}:

\begin{enumerate}
  \item at regular intervals, rank all stocks on the characteristic;
  \item group them into portfolios: deciles (ten groups), quintiles (five) or the 30/40/30 split used for HML;
  \item hold each portfolio for a period, and record its return.
\end{enumerate}

A \textbf{long--short} or \textbf{high-minus-low}\index{factor!HML (value)|textbf} portfolio buys the top group and sells the bottom one. If the characteristic is related to expected returns\index{return!expected}, the long--short portfolio\index{long--short portfolio} will earn a positive average return.

\subsection{Factors and factor models}

When a long--short portfolio captures a pattern shared by many stocks, it can serve as a factor in a model of expected returns, such as the Fama--French\index{Fama--French factors} three-factor model of Section~\ref{sec:03-factors} (Equation~\eqref{eq:03-ff3}). This chapter calls a stock's or fund's exposures to the factors its \textbf{loadings}\index{factor model!factor loadings|textbf}. Estimated by regression, they describe the fund's style, while its alpha relative to the model measures what the factors do not explain.

\textbf{Worked example.} A company has book equity of \$2 billion and 100 million shares trading at \$50, a market value of \$5 billion. Its book-to-market ratio is $2/5 = 0.4$ and its price-to-book ratio 2.5. If it had spent \$1 billion on research over the past five years that accounting treated as expenses, and that spending were instead counted as an asset, its adjusted book-to-market would be closer to $3/5 = 0.6$.

\providecolor{fenavy}{HTML}{1B2A6B}

\begin{figure}[!htbp]
  \centering
  \scalebox{0.8}{%
\begin{tikzpicture}[x=1cm, y=0.7cm]
    % Baseline
    \draw (-0.3,0) -- (6.9,0);

    % Market value: $5.0bn
    \fill[fenavy] (0,0) rectangle (1.4,5);
    \node[anchor=south] at (0.7,5) {\$5.0bn};

    % Book value as reported: $2.0bn
    \fill[gray!45] (2.6,0) rectangle (4.0,2);
    \node[anchor=south] at (3.3,2) {\$2.0bn};

    % Book value with R&D counted: $2.0bn + $1.0bn of R&D capital
    \fill[gray!45] (5.2,0) rectangle (6.6,2);
    \filldraw[fill=gray!20, draw=gray!60, dashed] (5.2,2) rectangle (6.6,3);
    \node at (5.9,2.5) {R\&D};
    \node[anchor=south] at (5.9,3) {\$3.0bn};

    % Bar labels
    \node[anchor=north, align=center] at (0.7,-0.15)
      {market value\\(price $\times$ shares)};
    \node[anchor=north, align=center] at (3.3,-0.15)
      {book value\\(accounting)};
    \node[anchor=north, align=center] at (5.9,-0.15)
      {book value with\\R\&D counted};

    % Ratios
    \node[anchor=north west, align=left] at (7.2,4.6)
      {book-to-market:\\$2/5 = 0.4$ as reported\\$3/5 = 0.6$ adjusted};
  \end{tikzpicture}%
}
\caption{Book value, market value and intangibles, from the worked example. Counting past research spending as an asset raises book value from \$2 billion to \$3 billion and book-to-market from 0.4 to 0.6, moving the firm toward the value end of a sort.}
\label{fig:08-book}
\end{figure}

\section{How it works}
\label{sec:08-how}

\textbf{Measuring cheapness.} A value stock is one whose price is low relative to a fundamental measure of the firm. The most studied measure is the book-to-market ratio. Other measures divide earnings, cash flow, dividends or sales by price. Each tells a slightly different story: book value measures past investment, earnings measure current profitability, dividends measure payouts.

\textbf{Sorting into portfolios.} The academic approach sorts all stocks on the chosen ratio and compares the returns of the cheapest group with the most expensive. Fama and French's \textbf{HML} (``high minus low'') factor does this carefully (Figure~\ref{fig:08-hml}):

\begin{enumerate}
  \item Each June, stocks are split into small and big by market value, and into three groups by book-to-market: growth (bottom 30\%), neutral and value (top 30\%).
  \item Six value-weighted portfolios result: small growth, small neutral, small value, big growth, big neutral, big value.
  \item HML is the average return of the two value portfolios minus the average return of the two growth portfolios. It is long cheap stocks and short expensive ones, with equal weight on small and big firms.
\end{enumerate}

\providecolor{fenavy}{HTML}{1B2A6B}

\begin{figure}[!htbp]
  \centering
  \scalebox{0.8}{%
\begin{tikzpicture}[
    cell/.style={minimum width=2.9cm, minimum height=1.15cm, inner sep=0pt,
                 draw=white, line width=1.2pt},
    growth/.style={cell, fill=gray!45},
    neutral/.style={cell, fill=gray!15},
    value/.style={cell, fill=fenavy, text=white},
  ]
    % Column headings: book-to-market groups
    \node[align=center] at (0,1.25)   {bottom 30\% B/M\\growth};
    \node[align=center] at (2.9,1.25) {middle 40\%\\neutral};
    \node[align=center] at (5.8,1.25) {top 30\% B/M\\value};

    % Row headings: size groups
    \node[anchor=east] at (-1.6,0)     {small};
    \node[anchor=east] at (-1.6,-1.15) {big};

    % The six portfolios
    \node[growth]  at (0,0)       {small growth};
    \node[neutral] at (2.9,0)     {small neutral};
    \node[value]   at (5.8,0)     {small value};
    \node[growth]  at (0,-1.15)   {big growth};
    \node[neutral] at (2.9,-1.15) {big neutral};
    \node[value]   at (5.8,-1.15) {big value};

    % HML definition
    \node at (2.9,-2.6)
      {HML $= \tfrac{1}{2}(\text{small value} + \text{big value})
             - \tfrac{1}{2}(\text{small growth} + \text{big growth})$};
    \node[text=gray!40!black, font=\small] at (2.9,-3.4)
      {rebuilt each June; held for a year};
  \end{tikzpicture}%
}
\caption{How HML is built. Stocks are sorted into two size groups and three book-to-market groups. HML buys the two value portfolios and sells the two growth portfolios, so that small and big firms count equally.}
\label{fig:08-hml}
\end{figure}

The portfolios are held for a year and rebuilt the next June, using book equity from the fiscal year ending in the previous calendar year so that the accounting data would have been public (Section~\ref{sec:02-lookahead}). This design keeps HML's turnover low, about 3\% a month (Chapter~\ref{ch:06}; Novy-Marx and Velikov 2016).

\textbf{Why value might earn more.} Two families of explanation compete:

\begin{itemize}
  \item \textbf{Risk.} Value firms are often distressed or inflexible: they have assets in place that are hard to cut back in recessions. If they do worst in bad economic times, investors demand higher expected returns to hold them (Fama and French 1993; Zhang 2005).
  \item \textbf{Mispricing.} Investors extrapolate past growth too far into the future, overpaying for glamorous growth firms and shunning unexciting value firms. As the expected growth fails to arrive, prices correct (Lakonishok, Shleifer and Vishny 1994).
\end{itemize}

The distinction matters for the future. A risk premium should persist, though it may arrive unevenly. A mispricing can shrink as investors learn about it and as capital flows in to exploit it.

\section{The claim}
\label{sec:08-claim}

Stated so that it can fail:

\begin{callout}
\textbf{Claim.} Stocks that are cheap relative to fundamentals (high book-to-market, earnings-to-price, or cash-flow-to-price) earn higher risk-adjusted average returns than expensive growth stocks out of sample and after transaction costs.
\end{callout}

The last clause is the key. The value premium is well documented in the data that discovered it. The question is whether it held afterward, in other markets, with other measures, and at a size an investor could hold. Three tests follow from the claim:

\begin{itemize}
  \item \textbf{Time.} Does HML earn a positive premium after the original sample of Fama and French (1992), which ended in 1990, and in particular after 2007?
  \item \textbf{Breadth.} Does the premium appear with earnings, cash flow and dividends as well as book value, among large stocks as well as small ones, and in other countries?
  \item \textbf{Implementability.} Does it survive costs, and could an investor live with its \textbf{drawdowns}, the falls from a previous peak, long enough to collect it?
\end{itemize}

\section{The original evidence and its interpretations}
\label{sec:08-evidence}

\textbf{Before the factor.} Basu (1977) found that stocks with high earnings-to-price ratios earned higher risk-adjusted returns than the capital asset pricing model predicted. Rosenberg, Reid and Lanstein (1985) found the same for book-to-market. These were treated as anomalies: puzzles for the CAPM rather than a new principle.

\textbf{Fama and French (1992, 1993).} Using U.S. stocks from 1963 to 1990, Fama and French (1992) showed that size and book-to-market together captured the \textbf{cross-section of average returns}, the differences in average return across stocks, and that market beta, the CAPM's only variable, added nothing. Fama and French (1993) turned the finding into the three-factor model, with HML as the value factor, and interpreted it as compensation for a risk related to financial distress. The papers changed the field: the three-factor model became the standard benchmark for judging investment performance for two decades.

\textbf{Lakonishok, Shleifer and Vishny (1994)} offered the competing interpretation. Value stocks did not appear riskier in the ways that mattered: they did not do notably worse in bad markets or recessions. Instead, investors seemed to extrapolate past growth: \textbf{glamour} firms with strong past growth were priced as if it would continue, and it did not.

\textbf{Breadth.} Fama and French (1998) found value premiums in many international markets, and Asness, Moskowitz and Pedersen (2013) found value effects across countries and asset classes, from stocks to bonds to currencies to commodities. By the 2010s, value was one of the best-documented patterns in finance.

\textbf{The drought.} From 2007 to 2020, U.S. value stocks trailed growth stocks by a large cumulative margin, the longest and deepest drawdown in HML's history. Explanations multiplied.

\begin{itemize}
  \item \textbf{Mismeasured book value.} Book value had become a poor measure of a firm's assets because spending on research, software and brands is expensed rather than recorded as an asset. Eisfeldt, Kim and Papanikolaou (2022) showed that adding \textbf{capitalized} intangible investment, recorded as an asset rather than an expense, to book equity produced a value factor that performed substantially better, including in recent decades.
  \item \textbf{Interest rates and crowding.} Others pointed to falling interest rates, which raise the value of distant growth cash flows, or to crowding by investors following the same strategy.
  \item \textbf{Noise.} Fama and French (2021) split 1963--2019 in half and found much lower value premiums in the second half, especially among large stocks, but showed that the volatility of returns made it impossible to reject either that the premium was unchanged or that it had fallen to zero.
  \item \textbf{Revaluation.} Arnott, Harvey, Kalesnik and Linnainmaa (2021) argued that much of the drought reflected value stocks becoming cheaper relative to growth stocks, a revaluation that cannot continue indefinitely, rather than a disappearing premium.
\end{itemize}

\section{Lab~\thechapter: replicating Fama and French (2021) and testing the value drought}
\label{sec:08-lab}

\begin{callout}
\textbf{Lab~\ref{ch:08} at a glance.} \emph{Claim on trial\index{claims-on-trial method}:} cheap stocks earn higher risk-adjusted returns than expensive growth stocks out of sample and after costs. \emph{Data:} Kenneth French's library: the Fama--French factors from July 1926; the six size and book-to-market portfolios; univariate portfolios formed on book-to-market, earnings-to-price, cash flow-to-price and dividend yield; and HML factors for North America, Europe, Japan, Asia Pacific ex Japan and the developed markets outside the United States, from 1990. Optionally, the Jensen--Kelly--Pedersen (2023) data extends the test to dozens of countries. \emph{Code:} repository folder \texttt{ch08-value/}, with a guided notebook.
\end{callout}

\textbf{The setup.} Lab~\ref{ch:08} replicates the empirical evaluation of Fama and French (1993, 2021) across four distinct historical eras and global markets:

\begin{itemize}
  \item \textbf{Part A, replication.} HML's average return, t-statistic, Sharpe ratio and CAPM alpha over the full sample and in four eras: 1926--1963, the original Fama--French sample of 1963--1990, 1991--2006, and 2007 to the present. A formal test asks whether the post-2007 premium differs from the earlier one by more than chance would allow.
  \item \textbf{Part B, stress tests.} Section~\ref{sec:08-stress} describes these.
  \item \textbf{Part C, other markets.} International HML before and after 2007.
  \item \textbf{Part D, living with value.} Drawdowns, the longest time spent below a previous peak, the worst ten-year period, and a bootstrap of ten-year outcomes used to size a value tilt (Section~\ref{sec:08-investors}).
\end{itemize}

\textbf{What readers will find.} Lab Figure~\ref{ch:08}.1 shows the growth of a dollar in HML over the century, with the eras marked. Lab Figure~\ref{ch:08}.2 sets out the premium in each era with its 95\% interval: the test of the claim over time is whether the post-2007 interval excludes the earlier premium, excludes zero, or neither. Fama and French's own split of 1963--2019 found the last of these, and readers can check whether the years since 2019 have changed the answer. Lab Figure~\ref{ch:08}.5 shows how long value can disappoint, and Lab Figure~\ref{ch:08}.6 turns that history into a range of ten-year outcomes.

\section{Stress tests}
\label{sec:08-stress}

\textbf{1. Other measures of cheapness.} If the value premium reflects something about cheap firms, it should appear whether cheapness is measured with book value, earnings, cash flow or dividends. If it appears only with book-to-market, it may be an artifact of accounting. The lab compares top-30\%-minus-bottom-30\% spreads for all four measures, over the full sample and before and after 2007. Book value is the measure most exposed to the intangibles problem, because research spending and brand building reduce book equity directly; earnings and cash flow are affected too, but less.

\textbf{2. Large stocks only.} Many anomalies are strongest among small stocks, which are expensive to trade and hold little of the market's value (Section~\ref{sec:02-microcaps}). The $2 \times 3$ portfolios separate value among big stocks from value among small stocks. Fama and French (2021) found that the decline in the second half of their sample was much larger for big stocks; the lab extends their comparison to the present.

\textbf{3. Intangible-adjusted book equity.} Eisfeldt, Kim and Papanikolaou (2022) capitalized research and development and a share of selling, general and administrative spending into book equity. Their intangible-adjusted value factor performed better throughout, including in the drought years. Replicating this requires firm-level accounting data that is not freely available, so the lab does not reproduce it; an exercise outlines how to do so with access to Compustat. Readers should note what the result implies: the drought was partly a measurement problem, not only a failure of the idea that cheap stocks earn more.

\textbf{4. After costs.} Value is a slow strategy. With turnover of about 3\% a month, a one-way trading cost of 20 to 50 basis points reduces its average return by roughly 0.3 to 0.7 percentage points a year (Equation~\eqref{eq:06-net}: $4 \times 0.03 \times 12 = 1.44$ times the one-way cost each year). Costs are not what threatens the value premium.

\textbf{5. Other countries.} If value works because of something general about how markets price cheap firms, it should appear outside the United States. The French international factors allow a test since 1990, before and after 2007. The Jensen--Kelly--Pedersen data, covering many more countries, also invites the multiple-testing warning that Chapter~\ref{ch:10} develops: with dozens of countries, some will show significant premiums by chance, and some will show none.

\section{Verdict}
\label{sec:08-verdict}

Based on the published evidence through the end of the 2010s drought, to be confirmed by readers' runs of Lab~\ref{ch:08} on data through 2025: \textbf{weakens.}

\begin{center}
\small
\begin{tabularx}{0.9\linewidth}{@{}>{\raggedright\arraybackslash}p{0.24\linewidth}>{\raggedright\arraybackslash}X>{\raggedright\arraybackslash}p{0.21\linewidth}@{}}
\toprule
Test & Finding & Effect on the claim \\
\midrule
Full century of U.S. data & HML earned a positive average premium, with a t-statistic well above conventional thresholds & Supports \\
After the original sample, and especially after 2007 & Much smaller, and for large U.S. stocks close to zero or negative over long stretches; statistically hard to distinguish either from the old premium or from zero & Weakens \\
Other measures of cheapness & The premium appears with earnings and cash flow as well as book value, though not identically & Supports breadth \\
Large stocks & Weaker than among small stocks, particularly in recent decades & Weakens \\
Intangible-adjusted book equity & A better-measured value factor performs better, including recently & Suggests part of the drought is mismeasurement \\
Other countries & Present in many markets, but uneven across countries and periods & Supports, with variation \\
Costs & Small relative to the premium & No effect \\
\bottomrule
\end{tabularx}
\end{center}

The claim that cheap stocks earn more than expensive ones has a century of support and a sound economic logic, but its size, its reliability and the measure of cheapness that captures it are all less certain than the field believed in the 1990s. The verdict is dated by the sample tested; a value premium that returned strongly for a decade would move it toward survives, and another decade like 2007--2020 toward fails.

\section{What investors can use}
\label{sec:08-investors}

The central lesson of value's history is not about its average return but about its patience requirement.

\textbf{Size the tilt you can hold through ten bad years.} A \textbf{value tilt} of weight $w$ adds $w$ times HML to a market portfolio, so its return relative to the market is $w$ times HML. Lab~\ref{ch:08}'s Part D resamples HML's history to estimate a bad-case ten-year outcome, the 5th percentile, and reports the largest tilt whose bad case stays within a shortfall you specify. An investor who would abandon a strategy after trailing the market by 10\% over a decade should hold no larger a tilt than that, because abandoning it after the losses and before the recovery is the worst outcome of all.

\textbf{Expect long droughts.} Even if the premium is real, the lab's drawdown\index{drawdown} statistics show how long HML can spend below a previous peak. A premium that arrives once a decade, irregularly, is a premium most investors will not stay to collect.

\textbf{Prefer measures that are harder to distort.} Composite measures that combine several value signals, and book equity adjusted for intangible investment, are less exposed to accounting changes than raw book-to-market.

\textbf{Questions to ask a value manager:}

\begin{enumerate}
  \item Which measure of cheapness do you use, and how did it perform from 2007 to 2020?
  \item How much of your historical return comes from small stocks you could not hold at your current size?
  \item How do you handle firms whose book value omits most of their assets, such as software, pharmaceutical and brand-driven companies?
  \item What is your worst ten-year relative performance in your backtest\index{backtest} and in live trading, and how did clients respond?
  \item Is your recent improvement a return of the premium, or value stocks becoming more expensive relative to growth?
\end{enumerate}

\begin{callout}
\textbf{What Failed: confident calls on value's death and its return.} As the drought deepened in the late 2010s, confident claims multiplied in both directions. Some declared value permanently dead, arbitraged away by the factor-investing industry or made obsolete by a technology-driven economy. Others declared, year after year, that a sharp comeback was imminent because value had become historically cheap relative to growth. Both claims ran ahead of the evidence. A premium whose year-to-year volatility\index{volatility} is large relative to its average cannot be declared dead after a decade of losses, as Fama and French (2021) showed formally; and cheapness relative to growth, while it raises expected returns, says little about when they will arrive. Value rebounded strongly in 2021 and 2022 and then lagged again. The claim that anyone could confidently time value's death or its revival \textbf{fails}.
\end{callout}

\section{Open questions}
\label{sec:08-open}

\begin{itemize}
  \item \textbf{Has intangible investment changed what ``cheap'' means?} If so, what is the right measure of a firm's invested capital, and does a properly measured value premium look stable?
  \item \textbf{Risk or mispricing?} The drought is a natural experiment\index{causal inference!natural experiment}. Did value's losses coincide with bad times for investors, as a risk story would require, or with good times?
  \item \textbf{Interest rates.} How much of value's performance since 2007 is explained by falling, then rising, interest rates acting on the long-duration\index{duration} cash flows of growth firms?
  \item \textbf{International divergence.} Why has the value premium been more persistent in some markets, such as Japan, than in the United States?
  \item \textbf{Crowding.}\index{crowding|textbf} Did the growth of value funds and factor products reduce the premium, and can crowding be measured in advance?
\end{itemize}

\section*{Exercises}
\addcontentsline{toc}{section}{Exercises}

\subsection*{Conceptual (all readers)}

\begin{enumerate}
  \item Explain the difference between the risk and mispricing explanations of the value premium. What evidence from the 2007--2020 drought would favor each?
  \item Why does HML use book equity from the previous fiscal year and rebalance only in June? What would go wrong with monthly rebalancing on the latest accounting data?
  \item A firm spends heavily on research and development. How does that spending affect its book-to-market ratio, and why might it be misclassified as a growth stock?
  \item Explain why a value premium of 4\% a year with 12\% annual volatility cannot be distinguished from zero over a 15-year period. How many years would be needed?
\end{enumerate}

\subsection*{Lab extensions (technical readers)}

\begin{enumerate}
  \setcounter{enumi}{4}
  \item Split the sample at 1990 instead of 2007. Which break looks more like a change in the premium?
  \item Compare value among large and small stocks after 2007. What does the comparison imply about who could actually have earned the premium?
  \item Run Part D with a tolerance of 5\% and 20\%. How does the largest tilt change, and what would you recommend to a cautious and to a patient investor?
  \item With the JKP data, count the countries in which the value premium has a t-statistic above 2. How many would you expect by chance?
  \item With access to firm-level accounting data, rebuild book equity by capitalizing research and development and part of selling, general and administrative spending, and rebuild HML. Does the post-2007 shortfall shrink?
  \item Measure the \textbf{value spread}, the gap in book-to-market between value and growth portfolios, over time. Does a wide spread predict higher subsequent HML returns? Apply Chapter~\ref{ch:05}'s out-of-sample tests.
\end{enumerate}